% !TEX TS-program = xelatex
% !TEX encoding = UTF-8 Unicode
\documentclass[11pt]{article}
\usepackage[margin=1in]{geometry}

\usepackage{xcolor}
\usepackage{hyperref}
\hypersetup{
  colorlinks=true,
  citecolor=violet
}

\usepackage[citestyle=authoryear, bibstyle=reading, entrykey=false, abstract=false, file=false]{biblatex}
\addbibresource{references.bib}

\usepackage{fontspec} 
\setmainfont{Charis SIL}


\title{Thesis Proposal}
\author{Wren Arbor}
\date{October 1st, 2025}

\begin{document}
\maketitle

\section{Thesis of the Thesis}
Aside from visual presentation of one's body, the sound of a person's voice is among the most strongly gendered aspects of behavior. Trans people know this quite intimately, as our speech is highly scrutinized by those who --  consciously or unconsciously -- seek to enforce the gender binary. In my thesis, I plan to investigate how this knowledge manifests in English-speaking transfeminine speakers' linguistic practice with respect to an aspect of pronunciation which is often highly associated with gender in western hegemonic gender normativity -- the pitch of the /s/ sound -- through a comparison of their speech in a conversation with a cis person versus in a conversation with a trans person.

\section{Brief Literature Review}

The above claim that trans speech is often more scrutinized by cisgender people for gender conformity, that trans people are aware of this, and that their knowledge may manifest in linguistic practice is based upon existing findings in the field of transgender linguistics. Investigating transgender knowledge, \textcite{heffernan_constructing_2023} finds that transfeminine speakers characterize a `cisgender listening subject' in their metalinguistic commentaries which casts negative judgements when linguistic gender norms are not followed, and that transfeminine speakers expressed varied ways of negotiating this, from accomodating these expectations by a change in linguistic behavior towards conformance to asserting gender to be self-determined, resisting these judgements, and talking as one would prefer.  \textcite{jessee_cisgender_2025} emperically confirm that transgender speech is indeed subject this greater gender scrutiny than cisgender speech, finding that cisgender listeners apply stricter gendered standards to speech when they are told the speaker is a transwoman than when told the speaker is a ciswoman. Although \textcite{parnell-mooney_transgender_2019} predates these more recent findings, he presents a design very similar to my own but with a focus on transmasculine and nonbinary speakers, in which it was found that these speakers do in fact pronounce /s/ with different pitch when talking to a trans person versus a cis person, but that variation was did not pattern across speakers.  


My choice to focus on transfeminine speakers is motivated by three reasons: the relative scarcity of sociolinguistic research which investigates transfeminine speech on the phonetic level, the emperical gap left by left by the exclusion of transfeminine speakers from \textcite{parnell-mooney_transgender_2019}, and the contributions to sociolinguistic theory, especially existing questions in the scholarship on /s/, that may arise by investigating transfeminine speech.  The scarcity of sociophonetic research on transfeminine speakers has been noted by \textcite{zimman_transgender_2020}; in a review of the place of trans people (and other groups which fall outside of the gender binary) in the language and sex, gender, and sexuality literature, they find that  in the Global South, research has primarily focused on third and fourth-gender groups  which bear similarities to transfeminine people, and in the Global north it has leaned towards transmasculine people, ``leaving trans feminine people in the Global North and trans masculine people in the Global South underrepresented.'' As I turn to the reasoning behind my choice to look at /s/, I will demonstrate this gap and discuss the ways in which a study of transfeminine interlocutor-oriented style-shifting (the sociolinguistics term for changes in an individual's speech depending on who they're talking with, see eg. \cite{rickford_addressee_1994}) could contribute new insights on existing theoretical questions raised by prior work on /s/.

Termed `sexuality's most iconic variable' by \textcite{calder_gay_2020}, /s/ has become one of the most well-studied sites of socially-relevant variation in speech (called `variables' in sociolinguistics) in the field of language and sex, gender, and sexuality. This scholarship has primarily oriented itself around finding differences in /s/ pronunciation and perception between sex, gender, and sexuality groups, individuals, and contexts, and attempting to explain those differences. The most commonly studied aspect of /s/ is frontness, an analouge of /s/ pitch; a fronter tongue position produces a higher-pitch (see \cite{fuchs_differences_2010}). Specifically, studies consistently find that women produce significantly fronter /s/ than men, and that gay men produce significantly fronter /s/ than straight men (eg. \cite{podesva_sexuality_2015, hazenberg_language_2012, flipsen_acoustic_1999}; also reviewed in \cite{calder_gay_2020}), and that /s/ is linked to women, gayness, and femininity by listeners (eg. \cite{campbell-kibler_intersecting_2011, maegaard_s_2015}).

Using the framework of indexicality, by which linguistic features come to index (`point to') social meanings, scholars have considered the social meaning of fronted /s/ to be femininity. That /s/ comes to be associated with gay masculinity is understood by \textcite{calder_sissy_2019} to be an example of a meaning of a higher indexical order -- a notion developed by \textcite{silverstein_indexical_2003} by which a linguistic feature comes to have new meanings through the process of adaptation to a new usage context. In this case, the use of fronted /s/ by gay men adapts its original meaning, femininity, to its new context, gayness. This structure of meanings comes to be through the imagining of gay men as more feminine, or more woman-like, than straight men. \textcite{calder_sissy_2019} sees this as an example of the ideological process \textcite{irvine_language_2000} term fractal recursivity by which a divison of society gets projected inside of itself-- gay men become ``men who are women'' through fractal recursivity, and thus, at their use of a `woman's' linguistic feature, new indexical order of /s/ emerges, going from a link from fronter /s/ to femininity to a link from fronter /s/ in men to gayness.

Although this framework of meanings works very well to describe the data on white, middle-class, cisgender westerners,  \textcite{calder_whose_2021}  points out that this structure, although `real' (that is, truly present and salient for the people it describes), does not work for those outside of this hegemonic group. If fronted /s/ were simply a feminine and gay-male index, trans people would be expected to pattern just the same as cis people. The fact that they do not (as seen in, eg. \cite{hazenberg_language_2012, podesva_sexuality_2015, zimman_hegemonic_2013}) also cannot be explained as simply another level of fractal recursivity; unlike sexuality identity, which can construe itself ideologically as a layer over and above gender and sex in the fractal recursive tree, this is untenable for transness, which is unavoidably present at all levels. Similarly, that these patterns do not hold outside of the sphere of whiteness (see eg. \cite{calder_whose_2022, jessee_cisgender_2025, maegaard_s_2015}) exposes the above understanding of the meaning structure of fronted /s/ as limited and incomplete. 

Using the concept of the indexical field, developed by \textcite{eckert_variation_2008}, which refers to the set of all meanings a linguistic feature can have, depending on context,  \textcite{calder_whose_2021} argues that the majority of the research on /s/ which does not fall into the more recent fields of raciolinguistics and trans linguistics, tacitly normalizes the white, cis, western, middle-class subset of the much more expansive indexical field. They argue that an account of how sex, gender, sexuality, (and race) are produced through language is incomplete without investigations of ``speaking and listening subjects outside of White and cisgender populations.'' (43) This further motivates my study of transfeminine speakers; I hope to find how they navigate this hegemonic meaning structure, and what meanings they use /s/ frontness to create. A more recent theorization of indexicality by \textcite{barrett_sexuality_2024} suggests that speakers and listeners bring their own indexical inventories to an interaction, and differences between these, or intentional exertion of power over another, can result in instances of misrecognition, by which the intended meaning of a n indexical behavior is misinterpreted. By looking at how transfeminine speakers navigate the possibility of misrecognition (eg. a fronted /s/ making them sound like a gay man, or a backed /s/ making them sound like a man rather than a lesbian), I hope evaluate this theoretical position.

Furthermore, I hope to raise questions about the status of `meaning' as the function of linguistic behavior. In the same way as literary theorists and feminists have attempted to move analysis in their fields away from the search for a meaning and towards an understanding of doing (\cite{deleuze_november_1987, butler_gender_1990}; see also \cite{livia_its_1997}), I feel it is possible that sociolinguistics needs to take this turn; my study of how transfeminine speakers deal with those who hope to impose meaning on their speech is an attempt to address this question: do they engage in the game of meaning-making, either by developing their own meanings, taking up up hegemonic meanings, or strategizing around the indexical inventory of their interlocutor, or do they reject `meaning' entirely, instead focusing on what their speech can do for them within whichever context they find themselves in. Are these theorizations even properly distinct?


Persuant to both the emperical gaps around transfeminine speech, and the theoretical questions just raised, I plan to investigate the \textbf{research question}: How do transfeminine speakers' pronunciations of English /s/ compare in speech targeted to a cisgender interlocutor versus in speech targeted to a transgender interlocutor with respect to frontness, and how does this vary across speakers of differing social backgrounds and identities? 

\section{Methodology Run-down}
I plan to investigate this by gathering data on /s/ pitch used in conversations between transfeminine speakers and a cis woman interlocutor, and transfeminine speakers and a trans woman interlocutor via acoustic analysis. I plan to analyse this data quantitavely using statistical methods, as well as qualitatively using information gathered in interviews between myself and the speakers, and, as time permits, discourse analysis.

\subsection{Participants}
I plan to investigate this question through the recruitment of 6 participants from the Reed College student body, 2 of which will be my  interlocutors -- a cis woman and a trans woman. The other 4 will be transfeminine people -- my speakers -- ideally diverse in their identification with respect to the binary gender `woman'.

\subsection{Procedure of Recording Meetings}
After recruitment, meetings will be scheduled such that each speaker has a short conversation (about 20min each) with the cis interlocutor, then trans interlocutor, then myself. Interlocutors will  additionally meet with me beforehand to tell me about their background and recieve orientation to what will be expected of them should they provide informed consent. 

Recording meetings with speakers will all take place in the Lab of Linguistics, and I will be present just before each conversation to orient participants to the experiment, confirm informed consent, and set up recording equipment. Prior to each of the conversations between the speakers and interlocutors, the structure of the conversation will be explained. The conversation will have two modules: first, the pair will be asked to collaboratively complete a crossword puzzle, and second they will be asked to spend the rest of the available time getting to know each other and discussing experiences or interests they find they have in common. The purpose of the first module is to act as an `icebreaker' to get the speakers talking, and to provide a conversational context in which the personal identities of the conversants are not forefronted. The purpose of the second module is to provide the opposite conversational context, in which personal identity is forefronted. This design is informed by \textcite{holmes-elliott_substance_2017} and \textcite{rickford_addressee_1994}, in which it was found that style-shifting depends on topic (in my design, task-oriented vs. identity oriented), and stance-taking (in my design, oppositional, [eg. disputing a crossword answer] vs relational [trying to find commonalities and relate to one another]). 

The recordings of the speakers within these meetings will be the subject of acoustic analysis. The recordings of participants (interlocutors and speakers) meetings with myself will be recorded also, mainly to be available to reference back to in favor of fallible memory or notes. In these meetings, I will get to know the participants, aiming to understand their sex, gender, and sexuality identities, especially the speakers' transgender journeys, backgrounds, experiences, and metalinguistic orientations. Within these interviews, participants and I would settle on a system of coding sex, gender, and sexuality identity, racial identity, as well other aspects of background (eg. years since transition, english monolingual vs. other languages, etc) which is comfortable for them, and these codings will be operationalized as variables the statistical analysis.


In acoustic analysis, I will measure /s/ pitch by the metric of center of gravity (COG), following the large majority of acoustic research on /s/ (see eg. \cite{fuchs_differences_2010, podesva_sexuality_2015, hazenberg_language_2012}). This would be achieved first through orthographic transcription of breath groups in Praat, forced alignment of segments in DARLA, and measurement of  /s/ token COG using a Praat script.

From there, I would quantitavely analyse the relationships between the identity and background variables (as decided through the interview process), whether the interlocutor was trans, and the /s/ realizations within the associated recording. I would use both descriptive and inferential statistical methods to understand the patterns in the data. Inferrential statistics would take the form of a linear mixed effects model in R.

 For the qualitative analysis, I would consider the information learned about each participant in my interviews with them, and apply this knowledge to form an understanding of their behavior with respect to /s/. Additionally, I may ask for consent to show certain transcribed portions of the speech for the purposes of discourse analysis.

\subsection{Funding}
Estimation and breakdown of participants' hours:
\begin{itemize}
\item 4 `Speaker' participants will each be present for one 1h15m meeting (setup + 3 approx. 20min recording sessions)

\item 2 `Interlocutor' participants will each be present for and one 30min meeting (for me to get to know them and orient them to the study) and four 35min meetings (setup + recording session with one `speaker' participant), totaling 2h50m per `interlocutor' participant. 

\item `Speaker' participants time would total to 4h45m, and `Interlocutor' participants time would total to 5h40m. 

\item \textbf{Participants' time would total 10h25m}.

\end{itemize}
\textbf{Setting a payrate of \$18 per hour, this would require an estimated total participants' compensation of \$187.50.  No further funding needs are expected,} as recording equipment will be provided by the Lab of Linguistics, and free software will be used for data collection and analysis.

\printbibliography[title=References with Annotations]
\end{document}